% !TEX root = ../main.tex
% 第一日第三讲（Yiannis）：相对 Langlands 作为 TQFT
\section{第一日第三讲（Yiannis）}

今天的第一场报告讲的是（4 维）TQFT，第二场讲的是把 Langlands 看作 TQFT。
这第三场讲的是把相对 Langlands 看作 TQFT。
这里，相对 = Langlands + 一些装饰。

先快速回顾前两场讲过的东西。
一切都是 4 维 TQFT：对一个 $d$ 维闭流形 $C$，配上一些态空间 $A_G(C)$，其范畴复杂度与维数相反，即 $A_G(C)$ 是一个 $4-d-1$-范畴，其中 $-1$ 是数、$0$ 是向量空间、$1$ 是范畴。

设 $G$ 是局部域或整体域上分裂的连通约化代数群。
下面是 Langlands 的 A 侧理论的三个典型例子，其中 $G$ 给出与 $C$ 相伴的 A 侧理论 $A_G(C)$。

\begin{exmp}
$C=$ 整体域 $F$，例如数域，或函数域 $\mathbb F_q(\gamma)$，其中 $\gamma$ 是某条光滑射影曲线。
此时 $C$ 是某种 3 维流形。
与之相应的 $A_G(C)$，在 $F$ 为数域时会是一些处处无分歧的自守函数 $\Fns(G(F)\backslash G(\mathbb A_F)/G(\hat O))$。

在函数域情形有两种理解方式。
一种与上面的数域情形相同。
另一种理解是：$\mathbb F_q$ 上曲线 $\gamma$ 的 $G$-丛的群胚的同构类集合。
\end{exmp}

\begin{exmp}
设 $F_v$ 是某个局部域，例如 $F_v=\mathbb Q_p$ 或 $\mathbb F_q((t))$。
它们是 2 维的。
于是我们配给它 $A_G(C)$，即 $G(F_v)$-表示的范畴。
\end{exmp}

\begin{exmp}
2 维的另一个例子是 $\overline{\mathbb F}_q(\gamma)$。
配给它的范畴是 $\Bun^\gamma_G$ 上的层范畴，记作 $\Shv(\Bun^\gamma_G)$。
\end{exmp}

在讨论上述例子的边界条件之前，先回顾前几讲里的一些配边例子。

设 $S$ 是域 $F$ 的一个有限位点集，$\mathcal O_S$ 是 $S$-整数环。
则 $C=\Spec\mathcal O_S$ 是一个带边流形，其边界为 $\bigsqcup_{v\in S}\Spec F_v$。
我们把 $C$ 看作从空集到边界 $\partial C$ 的配边；按 TQFT 的说法，它对应于一个从单位对象到边界上的理论 $A(\partial C)=\Rep G(F_v)$ 的映射。
给出这样的映射，等同于给出 $A(\partial C)$ 中的一个对象；我们取在 $S$ 中各点可能分歧的自守形式，即 $\Fns(G(F)\backslash G(\mathbb A_F)/G(\hat O_S))$。

上述例子展示了第一个例子的理论与第二个例子的理论之间的部分关系，告诉我们它们如何粘在一起。
第一个例子与第三个例子之间也有关系，由某种 Frobenius 迹给出。
下面给出这背后的几何图景。

\inlfig[0.36\textwidth]{f31-mappingcyl}

给定 $\mathbb F_q$ 上的光滑射影曲线 $\gamma$，有映射 $\gamma\to\Spec\mathbb F_q$。
如 David 所解释，$\Spec\mathbb F_q$ 应看作某种圆，因为其平展覆盖的 Galois 群本质上是 $\mathbb Z$。
圆上的一点对应于选取一个代数闭包（非典范的），即一个映射 $\Spec\overline{\mathbb F}_q\to\Spec\mathbb F_q$。
该点上的纤维，就是把曲线 $\gamma$ 基变换到所选代数闭包所得的 $\gamma_{\overline{\mathbb F}_q}$；在几何类比中（把代数闭包换成 $\mathbb C$），它是一个 Riemann 曲面。
一句话，$\gamma$ 应看作圆上的一族 Riemann 曲面。
但这族曲面在拓扑上并不平凡：存在一个 Frobenius 作用，也有 David 在第一讲里解释过的映射柱构造。
Frobenius 作用会告诉你绕圆一周的单值，这就让你把 $\gamma$ 看作某种 3 维的东西，并对该作用取「迹」。
存在一个由 Frobenius 作用给出的范畴迹，把第一个例子的理论与第三个例子的理论联系起来。

那么，什么是相对 Langlands？什么是带边 TQFT？
让我画一幅 David 没时间画的图\zy{原文此处说「让我画一幅 David 没时间画的图」，但 PDF 在这一句与下一页之间没有插入任何图；紧接着的正文又提到「图中右侧的平凡配边」。这里按原文照译，实际图形缺失。[译者按（无外部依据）：依扫描原样报告，未补绘。]}。
他明天会更多地讲这一点。

于是想法是这样的：在 David 的报告里，我们的范畴是某个带配边的流形。
对每个流形，都有一个典范但「无聊」的配边，从该流形到它自身（把流形与区间 $[0,1]$ 相乘）。
这只是恒等映射。
所以带边的理论允许你删掉一些边界，同时仍把它当作你的周期的合法输入。
这意味着，当你把范畴扩大时，可以删掉边界，并把平凡配边（图中右侧）看作从空集到该流形的态射。
每个带边流形都带有从空集出发的典范态射。
因此，无论这里的理论 $A_G(C)$ 是什么（配给流形 $C$ 的），它都会带有从单位元出发的典范映射。
所以新理论（带边的）更丰富：它不只给出某个范畴，还给出该范畴中的一个对象。
当然，这个「杰出元」需要更多数据才能产生，而在我们的理论里，这些数据就是某个带 $G$-作用的空间 $X$。
这个元素将记作 $\mathcal O_{G,X}$。

我想在不谈哈密顿空间与辛空间的前提下，谈论这些额外数据（「装饰」）。
玩具情形的原型是 $G$ 为有限群。
设 $X$ 是带 $G$-作用的向量空间。
此时（回忆 $\mathcal F_G(C)$ 是场论）
\[
  \mathcal F_G(C)=\Maps(C,\mathrm{pt}/G)=\Hom(\pi_1(C),G)/\sim
\]
其中 $\sim$ 是 $G$-共轭。
事实上 $\mathcal F_G(C)$ 有很多名字，如 $\Bun_G$ 与 $\Loc_G$（有限群情形下它们相同）。
与 $C$ 相伴的 A 侧理论是 $A_G(C)=\Fns(\mathcal F_G(C))$。
因此，边界条件应是 $\mathcal F_G(C)$ 上的一个杰出函数，我们把它记作 $\mathcal O_{G,X}$。
记号提示：它由 $G$-空间 $X$ 的额外数据给出。

取主 $G$-丛 $\mathcal E_G\to C$。
由于 $G$ 作用在 $X$ 上，我们可以构造 $C$ 上的一个向量丛，记作 $X\times^G\mathcal E_G$：取 $X$ 与 $\mathcal E_G$ 的积，再模掉对角作用。

于是，在 $\Bun_G=\mathcal F_G(C)$ 上可以构造空间
\[
  \Bun^X_G=\bigl\{(E_G,\sigma):\sigma=\text{section of }X\times^G E_G\bigr\}.
\]
与这组数据相伴的杰出函数 $\mathcal O_{G,X}$ 定义为
\[
  E_G\mapsto\#\{\sigma\}=\#(\Bun^X_G)_{E_G}=\#X^\phi.
\]
最后一个等式来自把 $E_G\in\Maps(C,\mathrm{pt}/G)$ 与单值 $\phi\in\Hom(\pi_1(C),G)/\sim$ 等同起来。
此时 $\#\{\sigma\}$ 不过是 $X$ 在该单值作用下的不动点个数。

然而在 Langlands 的背景下，$\mathcal O_{G,X}$ 的上述两种描述之间有差别，但哲学是一样的。

玩具情形意在让你初步了解空间 $X$ 如何产生边界条件。
其实 $X$ 不必是向量空间，它可以只是一个集合。

现在回到我们的三个例子：在这些情形中 $\mathcal O_{G,X}$ 是什么？
下表是一个总结。

\begin{center}
\scriptsize
\renewcommand{\arraystretch}{1.3}
\begin{tabular}{|p{0.135\textwidth}|p{0.235\textwidth}|p{0.215\textwidth}|p{0.175\textwidth}|}
\hline
 & 例 1 & 例 2 & 例 3\\\hline
$C$\newline[维数 $d$ 的流形]
 & 整体域 $F$，[$d=3$]，\newline 例如数域、函数域 $\mathbb F_q(\gamma)$
 & 局部域 $F_\nu$，[$d=2$]，\newline 例如 $\mathbb Q_p$、$\mathbb F_q((t))$
 & $\overline{\mathbb F}_q(\gamma)$，\newline [$d=2$]\\\hline
$A_G(C)$\newline[$(4-d-1)$-范畴]
 & $\Fns(G(F)\backslash G(\mathbb A_F)/G(\hat O))$，\newline（函数域情形可看作 $\Bun^\gamma_G$）
 & $G(F_\nu)$-表示
 & $\Shv(\Bun^\gamma_G)$\\\hline
$\mathcal O_{G,X}$
 & 周期函数 $\mathcal P_X$
 & $\Fns(X(F_\nu))$
 & 把常值层 $\mathbb C\in\Shv(\Bun^X_G)$ 前推到 $\Shv(\Bun_G)$\\\hline
$B_{\check G}(C)$
 & 落入 $\check G$ 中的 Langlands 参数上的函数
 & 上同调 Langlands 对偶层（局部 Langlands 参数）
 & $\QCoh^!(\Loc^\gamma_{\check G})$\\\hline
$\mathcal O_{\check G,\check X}$
 & 某些 L-函数
 & （将在后续讲义中出现）
 & （将在后续讲义中出现）\\\hline
\end{tabular}
\end{center}

先看第三个例子。
$\mathcal O_{G,X}$ 应当是空间 $\Bun_G=\Bun^\gamma_G$ 上的某个层。
$X$ 的额外数据给出 $\Bun_G$ 上的空间 $\Bun^X_G$。
把常值层前推，就得到 $\Shv(\Bun_G)$ 中的一个元素。

现在来看看另外两个例子中的边界条件是什么。
先想象一下 Frobenius 迹在这里会给出什么。
对第一种情形，我们应当给出某些自守函数。
在函数域情形，它可以看作曲线 $\gamma$ 的 $\Bun_G$ 上的函数。
这种情形下边界条件的函数将称作周期函数，记作 $\mathcal P_X$，即 $\mathcal O_{G,X}=\mathcal P_X$。
这里 $X$ 是 $\mathbb F_q$ 上带 $G$-作用的簇。
函数定义在 $G$-丛上，所以我们把 $G$-丛的记号改为 $[g]$（此前是 $\mathcal E_G$）。
改记号的原因是我们把它们与双商中的阿代尔点等同起来。
现在函数定义为
\[
  \mathcal P_X([g])=\#\{\text{sections of }[g]\times^G X\to\gamma\}
\]
（这个数有限吗？可以是有理数吗？）

而若把 $X$ 的 $\hat O$-点的示性函数记作 $\mathbf 1_X$，周期函数还有另一种写法：
\[
  \mathcal P_X([g])=\sum_{\gamma\in X(F)}\mathbf 1_X(\gamma g)
\]
其中 $F$ 是整体域。

这种自守函数的构造对更一般的一类函数也有效。
对 $X(\mathbb A)$ 上具有所需性质的函数 $\Phi$（比如速降、紧支集、在 $X(\hat O)$ 平移下不变），可以把 $\Phi$ 沿 $X(F)$ 中点的左平移求和。
因为当 $X$ 是仿射簇时 $X(F)$ 是离散集，这个和有定义。
所得函数在左 $G(F)$ 作用下不变，因为集合 $X(F)$ 在 $G(F)$ 下不变。
这类构造一般称作 theta 级数；在具体语境中也可能称作 Poincaré 级数、（伪）Eisenstein 级数等。

最后，第二个例子要给出的边界条件是什么？
我们要放的是某个 $G(F_v)$-表示。
所以它不过是 $X(F_v)$ 上的函数空间。
它是 $G(F_v)$-表示，因为 $G$ 作用在 $X$ 上。

现在花点时间解释 $\Bun^\gamma_G(\mathbb F_q)$ 与阿代尔点 $G(\mathbb F_q)\backslash G(\mathbb A)/G(\hat O)$ 是如何等同的。

给定 $G$-丛 $E$，可以找到曲线 $\gamma$ 的一个开子集 $U$，使 $E$ 限制到 $U$ 上同构于平凡 $G$-丛。
这个同构当然不是给定的，但可以选取。
$U$ 在 $\gamma$ 中的补集是有限个点。
所以要得到 $E$，只需把 $U$ 上的 $G$-丛与每个不在 $U$ 中的点周围形式圆盘上的 $G$-丛粘起来。
这个等同的细节留作练习。

以上所谈的只是 Langlands 的一侧，通常称作理论的 A 侧。
但 Langlands 有两侧，以及两侧之间的对偶。
那么 B 侧是什么？
B 侧应当是与对偶群 $\check G$ 相伴的理论。

回忆第一个例子里，理论的 A 侧只是函数 $\Fns(\Bun_G(F))$。
B 侧理论 $B_{\check G}(C)$ 也是函数，但粗略地说，是到 $\check G$ 中的 Langlands 参数上的函数。
这里的对偶是什么？
对偶是两个函数空间之间的同构，是 Fourier 变换的某种非交换版本。
对自守函数，你应当做谱分解，分解到 Hecke 算子的特征函数上去。
没有 Hecke 算子，一切都没有意义。
这些等价背后隐含地有 Hecke 算子，它们让你能对空间做谱分解，而在 B 侧它们应当是算子的联合特征值。
最非平凡的一点是：A 侧原本没有 Galois 理论，但 Hecke 算子把它与带 Langlands 参数的 B 侧联系起来，而 Langlands 参数与 Galois 表示相关。

在第三个例子里，A 侧理论是 $\Bun_G$ 上的某个层，B 侧理论是某个拟凝聚层 $\QCoh^!(\Loc^\gamma_{\check G})$。
关于它我不打算说什么，只是为了完整而提一句。

第二个例子里的 B 侧理论，对应于 A 侧（即 $G(F_v)$ 的表示），也应当是某种谱分解。
这个范畴 Langlands 对偶最著名的版本是由 Fargues 与 Scholze 证明的。
还有一个由朱歆文证明的版本。

相对 Langlands 对偶应当处理第三列，即边界条件。
对 B 理论也应当有相应的边界条件 $\mathcal O_{\check G,\check X}$。

比如在第一个例子里，边界条件是某些周期函数或 theta 理论，相应的 B 侧边界条件是某个「L-函数」。
我们关于这个对偶所有的证据，要么来自这种情形（我们对周期与 L-函数有所了解），要么来自第二种情形。
但这些我就不讲了。

关于由 $X$ 或 $\check X$ 给出的边界条件，我其实说了点谎。
事实上它们一般与各自余切空间的类比物相关联。
所以猜想的相对 Langlands 对偶，实际上是在 $(G,M)$ 与 $(\check G,\check M)$ 之间，其中 $M,\check M$ 是哈密顿 $G$-空间、$\check G$-空间。
在某些例子里，$M$ 可以解释为余切空间 $T^*X$。

最后给出一个数论的例子，把 theta 级数与周期等同起来；这就是 Tate 论文里的例子。
考虑 $X=\mathbb A^1$，带 $G=\mathbb G_m$ 的作用。
取函数
\[
  \Phi=\Bigl(\bigotimes_{\nu<\infty}\mathbf 1_{\mathbb Z_p}\Bigr)\otimes\Phi_\infty,
  \qquad \Phi_\infty\in\mathcal S(\mathbb R),
\]
其中可取 $\Phi_\infty=e^{-x^2}$（Gauss 函数）。
由这个 $\Phi$ 构造的 theta 级数是
\[
  \sum_{q\in X(\mathbb Q)}\Phi(qg),\qquad g=(1,\cdots,1,x)\in[G].
\]
现在，若一个有理数属于每个 $\mathbb Z_p$，它必须是整数。
于是
\[
  \sum_{q\in X(\mathbb Q)}\Phi(qg)=\sum_{n\in\mathbb Z}\Phi_\infty(nx)
  =\sum_{n\in\mathbb Z}e^{-n^2x^2}.
\]
